% vim:tw=80 ts=2 et sw=2 indentexpr= :
\documentclass[a4paper,12pt]{scrartcl}
\usepackage[utf8]{inputenc}
\usepackage[T1]{fontenc}
\usepackage[ngerman]{babel}
\usepackage{libertine} % kann man notfalls auch ignorieren, wenns nicht da ist
\usepackage{textcomp} % notfalls für €
\usepackage{hyperref}
\usepackage{stratum0doc}
\usepackage{longtable,booktabs}
\usepackage{graphicx}

\title{6.~Mitgliederversammlung des Stratum~0~e.~V.}
\date{6. Dezember 2015}

\begin{document}
\maketitle
%{\footnotesize\tableofcontents}

\hyphenation{Mit-glie-der-ver-samm-lung}

%%%%%%%%%%%
%% TOP 0 %%
%%%%%%%%%%%
\section{Organisatorischer Overhead}
\begin{description}
  \item[Zeit:] Sonntag, 6. Dezember 2015, 14:00
  \item[Ort:] Stratum 0 e.~V., Hamburger Straße 273a, Braunschweig
  \item[Anwesend:] 24 stimmberechtigte von 70 Mitgliedern (34\,\%), die Versammlung
    ist somit beschlussfähig.
  \item[Wahl des Protokollanten:] rohieb, ohne Gegenstimmen.
  \item[Wahl des Versammlungsleiters:] larsan, ohne Gegenstimmen.
	\item[Wahl der Wahlleitung:] Kah, ohne Gegenstimmen.
  \item[Veranstaltung eröffnet] durch den Versammlungsleiter um 14:27
  \item[Tagesordnung:] ohne Gegenstimmen angenommen
\end{description}


%%%%%%%%%%%
%% TOP 1 %%
%%%%%%%%%%%
\section{Jahresbericht und Entlastung des Vorstands}

\subsection{Finanzbericht}
chrissi\textasciicircum{} fasst die finanzielle Lage für das letzte Jahr
zusammen.  Der vollständige Finanzbericht kann im Wiki eingesehen
werden.\footnote{siehe
\url{https://stratum0.org/wiki/Datei:Finanzbericht_2015.pdf}}
Die Zahlen beziehen sich auf den Zeitraum von 1. Dezember 2014 bis 30. November
2015.

\paragraph{Überblick über die Finanzen im Jahr 2015:}
\begin{longtable}{|l|r|r|}
  \hline
  \textbf{Bereich} & \textbf{Einnahmen [€]} & \textbf{Ausgaben [€]} \\
  \hline
  \endfirsthead
  \multicolumn{3}{l}{\emph{(Fortsetzung von vorheriger Seite)}} \\
  \hline
  \textbf{Bereich} & \textbf{Einnahmen [€]} & \textbf{Ausgaben [€]} \\
  \hline
  \endhead
  \multicolumn{3}{l}{\emph{(Fortsetzung auf nächster Seite)}} \\
  \endfoot
  \endlastfoot
  Ideeller Bereich: Allgemein       & 27{.}242{,}50 & -12{.}130{,}44 \\
  \quad davon Mitgliedsbeiträge     & 12{.}990{,}00 &                \\
  \quad davon Spenden               &  4{.}951{,}66 &                \\
  \quad davon Veranstaltungsförderung &             &  -2{.}600{,}00 \\
  \quad davon Vereinsserver         &               &      -400{,}19 \\
  \hline
  Ideeller Bereich: Projekte        &  2{.}197{,}27 &  -1{.}284{,}19 \\
  \quad davon Bastelmaterial        &               &      -214{,}21 \\
  \quad davon Stickmaschine         &      156{,}12 &       -35{,}30 \\
  \quad davon Schneidplotter        &       87{,}90 &      -118{,}20 \\
  \quad davon 3D-Drucker            &       99{,}92 &       -40{,}90 \\
  \quad davon Freifunk              &  1{.}603{,}33 &      -875{,}58 \\
  \quad davon CoderDojo             &      250{,}00 &                \\
  \hline
  Ideeller Bereich: Space           &      286{,}15 & -15{.}914{,}46 \\
  \quad davon Miete und Nebenkosten &      146{,}15 & -13{.}368{,}27 \\
  \quad davon Einrichtung           &      140{,}00 &  -1{.}274{,}03 \\
  \quad davon Renovierung Space 2.0 &               &      -919{.}28 \\
  \hline
  Zweckbetriebe                     &        0{,}00 &         0{,}00 \\
  \hline
  Vermögensverwaltung               &        0{,}00 &         0{,}00 \\
  \hline
  Wirtschaftlicher Geschäftsbetrieb &  5{.}422{,}16 &  -5{.}027{,}94 \\
  \quad (Matekasse) && \\
  \hline
  Mankobuchungen                    &       39{,}66 &       -56{,}49 \\
  \hline
  Gesamt:                           & 35{.}187{,}74 & -34{.}413{,}52 \\
  \hline\hline
  \textbf{Gewinn/Verlust Gesamt:}   &      774{,}22 &                \\
  \hline
\end{longtable}

Größere Kostenpunkte dieses Jahr waren (neben laufenden Kosten) die anhaltende
Renovierung (Holodeck), der Zuschuss zum EasterHegg 2015, sowie als
durchlaufender Posten das Sponsoring der Volkswagen AG für unser CTF-Team zur
DEFCON 2014.  Die Mankobuchungen wurden zum Großteil durch die frei zugängliche
Verbrauchsmaterialkasse verursacht, hierfür soll in Zukunft ein Kassenbuch
angelegt werden. Die Einnahmen und Ausgaben im wirtschaftlichen Geschäftsbetrieb
(Getränkeverkauf) halten sich die Waage.

Die Gegenüberstellung der durchschnittlichen Einnahmen und laufenden
Verpflichtungen pro Monat\footnote{siehe auch
\url{https://stratum0.org/wiki/Verträge}} zeigt auch dieses Jahr noch ein
positives Bild:

\begin{longtable}{|lr|lr|} % tabular kann keine \footnotes
  \hline
  \multicolumn{2}{|c|}{\textbf{Einnahmen}} &
  \multicolumn{2}{|c|}{\textbf{Verpflichtungen}} \\
  \hline
  Mitgliedsbeiträge: & 1{.}002\,€ & Miete, Nebenkosten:      & 630\,€ \\
  Spenden:           &     380\,€ & Strom:                   & 240\,€ \\
                     &            & Internet:                &  42\,€ \\
                     &            & Server, Domain:          &  51\,€ \\
                     &            & Haftpflichtversicherung: &  12\,€ \\
  \hline\hline
  \textbf{Gesamt:} & \textbf{1{.}382\,€} & \textbf{Gesamt:} & \textbf{975\,€} \\
  \hline
\end{longtable}

Zu erwähnen ist allerdings, dass LINET sein regelmäßiges Sponsoring Mitte des
Jahres eingestellt hat und somit keine Firmensponsoren mehr existieren. Die
Mitgliederzahl wächst nicht mehr groß, in Summe handelt es sich um etwa 2 neue
Mitglieder. Die Mitgliederzahl sollte im Auge behalten werden, damit die
laufenden Kosten noch gedeckt werden können. Die Zahlung der Mitgliedsbeiträge
hat sich allerdings im Vergleich zum letzten Jahr stabilisiert. Die
Abschlagszahlung für Strom wurde von naturstrom um 20€ nach oben korrigiert, was
den steigenden Verbrauch im letzten Jahr widerspiegelt.

Insgesamt verfügt der Verein durchschnittlich über Mittel von etwa 6{.}900€, es
bestehen allerdings davon folgende Rücklagen:

\begin{center}
\begin{tabular}{|l|r|}
  \hline
  \textbf{Zweck der Rücklage} & \textbf{Betrag} \\
  \hline
  Erhöhung Mietsicherheit Space 2.0 &     160{,}00\,€ \\
  \qquad (wurde bisher von der Vermieterin nicht eingefordert) & \\
  Puffer Einnahmeausfälle           & 3{.}000{,}00\,€ \\
  \hline
\end{tabular}
\end{center}

Die Rückstellung für die Nachzahlung an naturstrom wurde aufgelöst, nachdem die
ausstehenden Forderungen geltend gemacht wurden. Als frei verfügbare Mittel
bleiben damit etwa 3{.}700€.

\paragraph{Ausblick 2016:} Es wird nach dem vollzogenen Umzug und Ausbau mit
einer Stabilisierung der Nebenkosten gerechnet.  Die  Nebenkostenabrechnungen
für 2012 und 2013 stehen seitens der Hausverwaltung weiterhin aus, der Vorstand
rechnet diesbezüglich aber nicht mit Nachzahlungen oder Gutschriften. Außerdem
dürfen wir weiterhin Zuwendungsbescheinigungen ausstellen.

Fragen aus der Runde:

\question{Inwiefern taucht das EasterHegg 2015 in der Buchhaltung auf?}
\answer{Bis auf den Punkt Veranstaltungsförderung (Saalmiete für
  CVJM-Hotel) gar nicht, wurde durch die ausgegliederte SMFW UG
  (haftungsbeschränkt) mit reneger als Geschäftsführer abgewickelt.}

\question{Besteht ein Plan, sich an der UG finanziell zu beteiligen?}
\answer{Der Plan besteht auf lange Sicht.}

\question{Gibt es einen Haushaltsplan für nächstes Jahr?}
\answer{Besteht nicht, außer dass Mitgliederzahlen bitte wieder
  steigen sollen. Ausgaben werden wie bisher erwartet.\footnote{unabhängig
  davon besteht eine Auflistung der laufenden Verpflichtungen, siehe oben.}}

\subsection{Rechenschaftsbericht}
Kasalehlia fasst für den gesamten Vorstand dessen Tätigkeiten im letzten Jahr
zusammen. Seine Präsentation kann im Wiki heruntergeladen
werden.\footnote{siehe
\url{https://stratum0.org/wiki/Datei:Rechenschaftsbericht_2015.pdf}}

\subsection{Bericht der Kassenprüfer}
Eine der Kassenprüfer hat den Verein inzwischen verlassen, da in der
Zwischenzeit keine andere Person beauftragt wurde, hat shoragan die Kasse
alleine geprüft. Hier stellt sich die Frage, ob man sicherheitshalber mehr als
zwei Kassenprüfer wählen sollte. Satzung, §8 Abs.~4 spricht von "`zwei
Rechnungsprüfern"', dies ist wohl als "`genau zwei"' zu interpretieren.

Die Kassenprüfung hat festgestellt, dass für eine Rechnung zu viel überwiesen
wurde, dies befindet sich jedoch schon im Prozess der Reperatur. Durch die
vielen Mankobuchungen in der Verbrauchsmaterialkasse spricht sich auch shoragan
dringend für die Einführung eines Kassenbuchs aus. Außerdem bittet er, darauf zu
achten, dass bei Einzahlung in die Matekasse (Backbone) möglichst nicht viele, kleine
Buchungen erstellt werden, sondern lieber mehr Geld am Stück eingezahlt wird, um
die Fehlerquote und auch die Arbeit zur Digitialisierung gering zu halten.
Außerdem hebt er hervor, dass auf Kassenzetteln in der Verbrauchsmaterialkasse
immer ersichtlich sein muss,
welche Dinge für den Verein bezahlt wurde und welche privat. Notfalls sollte man
sich beim Einkauf eine Zwischensumme auf dem Kassenzettel erstellen lassen.

Da die Mankobuchungen angesichts der Umsatzmenge keine gravierenden Fehler
darstellen, empfiehlt shoragan die Entlastung des Schatzmeisters.

\subsection{Entlastung des Vorstands}
\consensus{Vorstand entlastet}
larsan beantragt, den Vorstand in seiner Gesamtheit zu entlasten. Eine
Einzelentlastung wird nicht gewünscht. Der Vorstand enthält sich bei der
Abstimmung. Die Handabstimmung ergibt ohne Gegenstimmen, dass der Vorstand
entlastet werden soll.

\emph{[ 15:35--15:47: kurze Pause. ]}

\subsection{Jahresbericht}
Nach der Pause gibt larsan einen Überblick, was generell so im letzten Jahr aus
Sicht des Vereins passiert ist. Seine Präsentation kann im Wiki heruntergeladen
werden.\footnote{siehe
\url{https://stratum0.org/wiki/Datei:Jahresbericht2015.pdf}}

\emph{[ 16:15: ein Mitglied verlässt die Versammlung, dessen unausgefüllter
Wahlzettel wird vom Schatzmeister ungültig gemacht. ]}



%%%%%%%%%%%
%% TOP 2 %%
%%%%%%%%%%%
\section{Wahlen}
\emph{[ 16:39: ein stimmberechtigtes Mitglied wird nachträglich akkreditiert ]}

Als Wahlmodus wird wie bisher vorgeschlagen:
\begin{itemize}
  \item alle Ämter werden gleichzeitig auf einem Zettel gewählt (wenig Papier).
  Da die Kandidaten vorher noch nicht feststanden, wurden Stimmzettel mit
  Symbolen vorgedruckt und es wird jedem Kandidat ein Symbol zugewiesen.
  \item Zustimmungswahl (alles ankreuzen, was gewählt werden soll)
  \item Nur Kandidaten mit mindestens 50\% der Stimmen sind gewählt
  \item Besetzung der Ämter in der Reihenfolge Vorstandsvorsitzender, stellv.
    Vorsitzender, Schatzmeister, Beisitzer (nach Stimmenanzahl absteigend),
    Rechnungsprüfer.
\end{itemize}

Die Mehrheit der Versammlung stimmt per Handzeichen für diesen Wahlmodus. Als
Wahlhelfer stellen sich mjh, Emantor zur Verfügung, dagegen gibt es keine
Einwände.

Die Kandidaten stellen sich jeweils kurz vor. Die Wahlurne wird ordnungsgemäß
ausgekippt und mit Klebeband versiegelt.

Der Wahlgang wird eröffnet um 16:56, die Wahlabgabe wird um 16:58 geschlossen.
Die Wahlleitung zieht sich mit den Wahlhelfern zur Auszählung zurück und es gibt
eine kurze Pause.

\emph{[ 2 Mitglieder verlassen in der Pause die Versammlung ]}

Das Wahlergebnis wird um 17:22 verkündet. Leider wurde die Anzahl der
abgegebenen Stimmen nicht bekanntgeben. Da jedoch 25 Stimmzettel bei der
Akkreditierung ausgegeben wurden, von denen einer ungülgtig gemacht wurde, und
da bei der Wahl mehrmals 24 Stimmen auf einen Kandidaten entfielen, ist sehr
wahrscheinlich, dass 24 gültige Stimmzettel abgegeben worden sind.

Wahlergebnis:

\subsection*{Vorstandsvorsitzender}
\begin{tabular}{|c|l|l|l|}
  \hline
  \textbf{Symbol} & \textbf{Kandidat} & \textbf{Stimmen} & \textbf{Prozent} \\
  \hline
  \includegraphics[height=20pt]{images/mv20151206-wahl-1v-1.pdf} &
  rohieb (Roland Hieber) & 13 & 51{,}4\% \\
  \includegraphics[height=20pt]{images/mv20151206-wahl-1v-2.pdf} &
  larsan (Lars Andresen) & 23 & 95{,}8\% \\
  \hline
\end{tabular}

\elected{Vorstands\-vorsitzender}{larsan}{23}{24}
larsan nimmt die Wahl zum Vorstandsvorsitzenden an.

\subsection*{Stellv. Vorsitzender}
\begin{tabular}{|c|l|l|l|}
  \hline
  \textbf{Symbol} & \textbf{Kandidat} & \textbf{Stimmen} & \textbf{Prozent} \\
  \hline
  \includegraphics[height=12pt]{images/mv20151206-wahl-2v-1.pdf} &
  Kasalehlia (Hilko Boekhoff)  & 23 & 95{,}8\% \\
  \includegraphics[height=12pt]{images/mv20151206-wahl-2v-2.pdf} &
  larsan (Lars Andresen)       & 21 & 87{,}5\% \\
  \includegraphics[height=12pt]{images/mv20151206-wahl-2v-3.pdf} &
  rohieb (Roland Hieber)       & 20 & 83{,}3\% \\
  \hline
\end{tabular}

\elected{stellv. Vorsitzender}{rohieb}{20}{24}
Kasalehlia nimmt die Wahl zum stellv. Vorsitzenden nicht an. larsan ist schon
als Vorstandsvorsitzender gewählt und scheidet damit aus. rohieb nimmt die Wahl
an.

\subsection*{Schatzmeister}
\begin{tabular}{|c|l|l|l|}
  \hline
  \textbf{Symbol} & \textbf{Kandidat} & \textbf{Stimmen} & \textbf{Prozent} \\
  \hline
  \includegraphics[height=12pt]{images/mv20151206-wahl-schatzmeister-1.pdf} &
  chrissi\textasciicircum{} (Chris Fiege) & 24 & 100\% \\
  \hline
\end{tabular}

\elected{Schatzmeister}{chrissi\textasciicircum}{24}{24}
chrissi\textasciicircum{} nimmt die Wahl zum Schatzmeister an.

\subsection*{Beisitzer}
\begin{tabular}{|c|l|l|l|}
  \hline
  \textbf{Symbol} & \textbf{Kandidat} & \textbf{Stimmen} & \textbf{Prozent} \\
  \hline
  \includegraphics[height=12pt]{images/mv20151206-wahl-beisitzer-1.pdf} &
  Kasalehlia (Hilko Boekhoff)      & 24 & 100\% \\
  \includegraphics[height=12pt]{images/mv20151206-wahl-beisitzer-2.pdf} &
  larsan (Lars Andresen)           & 22 & 91{,}6\% \\
  \includegraphics[height=12pt]{images/mv20151206-wahl-beisitzer-3.pdf} &
  rohieb (Roland Hieber)           & 19 & 79{,}1\% \\
  \includegraphics[height=12pt]{images/mv20151206-wahl-beisitzer-4.pdf} &
  hanhaiwen (Helga Hansen)         & 21 & 87{,}5\% \\
  \includegraphics[height=12pt]{images/mv20151206-wahl-beisitzer-5.pdf} &
  DooMMasteR (Steffen Arntz)       & 16 & 66{,}6\% \\
  \includegraphics[height=12pt]{images/mv20151206-wahl-beisitzer-6.pdf} &
  tnias                            & 13 & 54{,}1\% \\
  \includegraphics[height=12pt]{images/mv20151206-wahl-beisitzer-7.pdf} &
  lichtfeind                       & 11 & 45{,}8\% \\
  \hline
\end{tabular}

\elected{Beisitzer}{Kasa\-lehlia}{24}{24}
\elected{Beisitzerin}{hanhaiwen}{21}{24}
\elected{Beisitzer}{DooMMasteR}{16}{24}
Kasalehlia nimmt die Wahl zum Beisitzer an. larsan und rohieb sind schon
gewählt und scheiden somit aus. hanhaiwen und DooMMasteR nehmen beide die Wahl
zur Beisitzerin bzw. zum Beisitzer an.

\subsection*{Rechnungsprüfer}
\begin{tabular}{|c|l|l|l|}
  \hline
  \textbf{Symbol} & \textbf{Kandidat} & \textbf{Stimmen} & \textbf{Prozent} \\
  \hline
  \includegraphics[height=12pt]{images/mv20151206-wahl-rechnungspruefer-1.pdf} &
  rohieb   & 13 & 54.1\% \\
  \includegraphics[height=12pt]{images/mv20151206-wahl-rechnungspruefer-2.pdf} &
  shoragan & 24 & 100\%  \\
  \includegraphics[height=12pt]{images/mv20151206-wahl-rechnungspruefer-3.pdf} &
  Angela   & 23 & 95.8\% \\
  \hline
\end{tabular}

\elected{Rechnungsprüfer}{shoragan}{24}{24}
\elected{Rechnungsprüfer}{Angela}{23}{24}
rohieb ist schon als stellv. Vorsitzender gewählt. shoragan und Angela nehmen
die Wahl zu Rechnungsprüfern an.


%%%%%%%%%%%
%% TOP 3 %%
%%%%%%%%%%%
\section{Satzungsänderungen}
\emph{[ 17:29: Durchzählen zur Kontrolle: 22 stimmberechtigte Mitglieder
anwesend. ]}

Satzungsänderungsanträge müssen laut Satzung §9 mindestens drei Viertel der
Stimmen erhalten, dies entspricht bei 22 anwesenden, stimmberechtigten
Mitgliedern mindestens 17 Stimmen.

Die folgenden Änderungsanträge zur Satzung wurden gestellt und sind hier
mit altem und neuem Wortlaut wiedergegeben (Unterschiede jeweils
kursiv):

\subsection{Änderungsantrag 1: Quoren für Fördermitglieder einschränken}
Die Fördermitgliedschaft war von Anfang an für Mitglieder gedacht, die den
Verein vorrangig finanziell unterstützen wollen und an Abstimmungen auf
Mitgliederversammlungen kein Interesse haben (s. Satzung, §4 (Mitgliedschaft),
Abs. 2). Leider wurde bei der entsprechenden Satzungsänderung auf der
Mitgliederversammlung 2013-12-07, die die Fördermitgliedschaft
einführte,\footnote{siehe \url{
https://stratum0.org/wiki/Datei:Mitgliederversammlung_2013-12-07.pdf?page=5}}
einige Formulierungen nicht angepasst. So zählen Fördermitglieder bisher noch
zum Quorum der Beschlussfähigkeit bei Mitgliederversammlungen, und in andere
Quoren, bei denen sie schon jetzt gar kein Stimmrecht haben.

Da nur stimmberechtigte Mitglieder abstimmen können, sollen die Quoren für
Beschlussfassung entsprechend anhand der anwesenden stimmberechtigten Mitglieder
berechnet werden:

\subsubsection*{§7 (Mitgliederversammlung), Abs. 2}
\begin{longtable}[c]{@{}ll@{}}
\toprule
\begin{minipage}[b]{0.47\columnwidth}\raggedright\strut
alter Text:
\strut\end{minipage} &
\begin{minipage}[b]{0.47\columnwidth}\raggedright\strut
neuer Text:
\strut\end{minipage}\tabularnewline
\midrule
\endhead
\begin{minipage}[t]{0.47\columnwidth}\raggedright\strut
Eine Mitgliederversammlung ist außerdem einzuberufen, wenn es das
Vereinsinteresse erfordert, oder wenn die Einberufung von mindestens
23\% der Vereinsmitglieder in Textform und unter Angabe des Zweckes und
der Gründe verlangt wird.
\strut\end{minipage} &
\begin{minipage}[t]{0.47\columnwidth}\raggedright\strut
Eine Mitgliederversammlung ist außerdem einzuberufen, wenn es das
Vereinsinteresse erfordert, oder wenn die Einberufung von mindestens
23\% der \emph{ordentlichen} Vereinsmitglieder in Textform und unter Angabe des
Zweckes und der Gründe verlangt wird.
\strut\end{minipage}\tabularnewline
\bottomrule
\end{longtable}

\subsubsection*{§7 (Mitgliederversammlung), Abs. 6}
\begin{longtable}[c]{@{}ll@{}}
\toprule
\begin{minipage}[b]{0.47\columnwidth}\raggedright\strut
alter Text:
\strut\end{minipage} &
\begin{minipage}[b]{0.47\columnwidth}\raggedright\strut
neuer Text:
\strut\end{minipage}\tabularnewline
\midrule
\endhead
\begin{minipage}[t]{0.47\columnwidth}\raggedright\strut
Jede satzungsmäßig einberufene Mitgliederversammlung wird als
beschlussfähig anerkannt, sofern mindestens 23\% der Mitglieder anwesend
sind. Falls dieser geforderte Anteil nicht erreicht wird, ist die darauf
folgende Mitgliederversammlung unabhängig von der Anzahl der erschienen
Mitglieder beschlussfähig. Auf diesen Umstand muss in der Einladung zur
Mitgliederversammlung besonders hingewiesen werden. Jedes ordentliche
Mitglied hat eine Stimme. Fördermitglieder sind berechtigt, an den
Versammlungen ohne Stimmrecht teilzunehmen.
\strut\end{minipage} &
\begin{minipage}[t]{0.47\columnwidth}\raggedright\strut
Jede satzungsmäßig einberufene Mitgliederversammlung wird als
beschlussfähig anerkannt, sofern mindestens 23\% der \emph{ordentlichen}
Mitglieder anwesend sind. Falls dieser geforderte Anteil nicht erreicht
wird, ist die darauf folgende Mitgliederversammlung unabhängig von der
Anzahl der erschienen Mitglieder beschlussfähig. Auf diesen Umstand muss
in der Einladung zur Mitgliederversammlung besonders hingewiesen werden.
Jedes ordentliche Mitglied hat eine Stimme. Fördermitglieder sind
berechtigt, an den Versammlungen ohne Stimmrecht teilzunehmen.
\strut\end{minipage}\tabularnewline
\bottomrule
\end{longtable}

\subsubsection*{§7 (Mitgliederversammlung), Abs. 7}
\begin{longtable}[c]{@{}ll@{}}
\toprule
\begin{minipage}[b]{0.47\columnwidth}\raggedright\strut
alter Text:
\strut\end{minipage} &
\begin{minipage}[b]{0.47\columnwidth}\raggedright\strut
neuer Text:
\strut\end{minipage}\tabularnewline
\midrule
\endhead
\begin{minipage}[t]{0.47\columnwidth}\raggedright\strut
Die Mitgliederversammlung fasst ihre Beschlüsse mit einfacher Mehrheit
der anwesenden Mitglieder, sofern in dieser Satzung nicht anders
geregelt. Bei Stimmengleichheit gilt ein Antrag als abgelehnt.
\strut\end{minipage} &
\begin{minipage}[t]{0.47\columnwidth}\raggedright\strut
Die Mitgliederversammlung fasst ihre Beschlüsse mit einfacher Mehrheit
der anwesenden \emph{stimmberechtigten} Mitglieder, sofern in dieser Satzung
nicht anders geregelt. Bei Stimmengleichheit gilt ein Antrag als
abgelehnt.
\strut\end{minipage}\tabularnewline
\bottomrule
\end{longtable}

\newpage
\subsubsection*{§9 (Satzungsänderungen), Abs. 1}
\begin{longtable}[c]{@{}ll@{}}
\toprule
\begin{minipage}[b]{0.47\columnwidth}\raggedright\strut
alter Text:
\strut\end{minipage} &
\begin{minipage}[b]{0.47\columnwidth}\raggedright\strut
neuer Text:
\strut\end{minipage}\tabularnewline
\midrule
\endhead
\begin{minipage}[t]{0.47\columnwidth}\raggedright\strut
Für Satzungsänderungen ist eine Dreiviertel-Mehrheit der anwesenden
Vereinsmitglieder erforderlich. Über Satzungsänderungen kann in der
Mitgliederversammlung nur abgestimmt werden, wenn auf diesen
Tagesordnungspunkt bereits in der Einladung zur Mitgliederversammlung
hingewiesen wurde und der Einladung sowohl der bisherige als auch der
vorgesehene neue Satzungstext beigefügt wurde.
\strut\end{minipage} &
\begin{minipage}[t]{0.47\columnwidth}\raggedright\strut
Für Satzungsänderungen ist eine Dreiviertel-Mehrheit der anwesenden
\emph{stimmberechtigten} Vereinsmitglieder erforderlich. Über
Satzungsänderungen kann in der Mitgliederversammlung nur abgestimmt
werden, wenn auf diesen Tagesordnungspunkt bereits in der Einladung zur
Mitgliederversammlung hingewiesen wurde und der Einladung sowohl der
bisherige als auch der vorgesehene neue Satzungstext beigefügt wurde.
\strut\end{minipage}\tabularnewline
\bottomrule
\end{longtable}

Zum Stimmrecht bei Änderung des Vereinszwecks gibt es mit
Änderungsantrag 2 und 3 zwei Varianten, die nach Annahme dieses Antrags
abgestimmt werden sollen.

\vote{Quoren für För\-der\-mit\-glieder ein\-schränken}{22}{0}{0}
Der Antrag wird per Handzeichen abgestimmt und einstimmig angenommen.


\subsection{Änderungsantrag 2: Stimmrecht bei Änderung des Vereinszwecks
auf ordentliche Mitglieder einschränken}
Dieser Antrag soll nur nach Annahme von Änderungsantrag 1 abgestimmt
werden und steht in Konkurrenz zu Änderungsantrag 3. Analog zu den
übrigen Bestimmungen sollen nur ordentliche Mitglieder Stimmrecht bei
Änderung des Vereinszwecks erhalten.

\subsubsection*{§9 (Satzungsänderungen), Abs. 2}
\begin{longtable}[c]{@{}ll@{}}
\toprule
\begin{minipage}[b]{0.47\columnwidth}\raggedright\strut
alter Text:
\strut\end{minipage} &
\begin{minipage}[b]{0.47\columnwidth}\raggedright\strut
neuer Text:
\strut\end{minipage}\tabularnewline
\midrule
\endhead
\begin{minipage}[t]{0.47\columnwidth}\raggedright\strut
Zur Änderung des Vereinszwecks ist die Zweidrittel-Mehrheit aller
Vereinsmitglieder erforderlich, wobei die Zustimmung der nicht
anwesenden Mitglieder per Textform erfolgen kann.
\strut\end{minipage} &
\begin{minipage}[t]{0.47\columnwidth}\raggedright\strut
Zur Änderung des Vereinszwecks ist die Zweidrittel-Mehrheit aller
\emph{ordentlichen} Vereinsmitglieder erforderlich, wobei die Zustimmung der
nicht anwesenden Mitglieder per Textform erfolgen kann.
\strut\end{minipage}\tabularnewline
\bottomrule
\end{longtable}

\vote{Stimmrecht bei Änderung des Vereinszwecks auf ordentliche Mitglieder
einschränken}{16}{3}{3} Die Abstimmung per Handzeichen ergibt 16 pro (72.7\%), 3
contra, 3 Enthaltungen. Der Antrag ist damit abgewiesen.


\subsection{Änderungsantrag 3: Fördermitglieder haben je eine Stimme bei
Änderung des Vereinszwecks}
Dieser Antrag soll nur nach Annahme von Antrag 1 abgestimmt werden und
steht in Konkurrenz zu Änderungsantrag 2. Da Fördermitglieder vorrangig
die gemeinnützigen Zwecke des Vereins fördern, sollen sie auch ein
Mitspracherecht bei deren Änderung haben.

\newpage
\subsubsection*{§9 (Satzungsänderungen), Abs. 2}
\begin{longtable}[c]{@{}ll@{}}
\toprule
\begin{minipage}[b]{0.47\columnwidth}\raggedright\strut
alter Text:
\strut\end{minipage} &
\begin{minipage}[b]{0.47\columnwidth}\raggedright\strut
neuer Text:
\strut\end{minipage}\tabularnewline
\midrule
\endhead
\begin{minipage}[t]{0.47\columnwidth}\raggedright\strut
Zur Änderung des Vereinszwecks ist die Zweidrittel-Mehrheit aller
Vereinsmitglieder erforderlich, wobei die Zustimmung der nicht
anwesenden Mitglieder per Textform erfolgen kann.
\strut\end{minipage} &
\begin{minipage}[t]{0.47\columnwidth}\raggedright\strut
Zur Änderung des Vereinszwecks ist die Zweidrittel-Mehrheit aller
Vereinsmitglieder erforderlich, wobei die Zustimmung der nicht
anwesenden Mitglieder per Textform erfolgen kann. \emph{Fördermitglieder haben
hierbei, abweichend von den übrigen Bestimmungen dieser Satzung, je eine
Stimme.}
\strut\end{minipage}\tabularnewline
\bottomrule
\end{longtable}

\withdrawn
Nach Diskussion im Plenum wird dieser Antrag wird vom Antragsteller
zurückgezogen. Argumente: es handelt sich hierbei nur um "`nettes"' und
"`demokratisches"' Verhalten gegenüber Fördermitgliedern, dieser können aber ja
trotzdem immer noch austreten, falls ihnen der Satzungszweck nicht mehr passt.
Außerdem könnten viele Fördermitglieder einige wenige ordentliche Mitglieder
überstimmen, was nicht gewünscht wird. Falls außerdem zusätzlich Änderungsantrag
4 "`Wechsel der Art der Mitgliedschaft vereinfachen"' angenommen würde (siehe
\ref{uxe4nderungsantrag-4-wechsel-der-art-der-mitgliedschaft-vereinfachen}),
könnten Fördermitglieder vor einer Zweckänderung zudem einfach wechseln.


\subsection{Änderungsantrag 4: Wechsel der Art der Mitgliedschaft vereinfachen}
\label{uxe4nderungsantrag-4-wechsel-der-art-der-mitgliedschaft-vereinfachen}
Da die Satzung bisher keinen Wechsel der Art der Mitgliedschaft vorsieht
(z.~B. von ordentlicher zu Fördermitgliedschaft), wird dieser Vorgang
vorstandsintern bisher als Aus- und Wiedereintritt behandelt, was (wegen
des Wiedereintritts) zwingend eine Vorstandsentscheidung mit Abstimmung
voraussetzt. Um den Aufwand für den Vorstand zu vereinfachen, kann diese
Abstimmung im Prinzip entfallen, da das Mitglied ja weiterhin Mitglied
bleibt. Es gibt hier für den Vorstand auch keinen sinnvollen Grund,
einen solchen Antrag abzulehnen -- das Mitglied sollte selbst am besten
wissen, wie viel Geld es am besten in den Verein investieren kann.

\subsubsection*{§4 (Mitgliedschaft), Abs. 5}
\begin{longtable}[c]{@{}ll@{}}
\toprule
\begin{minipage}[b]{0.47\columnwidth}\raggedright\strut
alter Text:
\strut\end{minipage} &
\begin{minipage}[b]{0.47\columnwidth}\raggedright\strut
neuer Text:
\strut\end{minipage}\tabularnewline
\midrule
\endhead
\begin{minipage}[t]{0.47\columnwidth}\raggedright\strut
Ein Austritt ist jederzeit möglich und wird durch Willenserklärung in
Textform gegenüber dem Vorstand vollzogen.
\strut\end{minipage} &
\begin{minipage}[t]{0.47\columnwidth}\raggedright\strut
Ein Austritt \emph{oder ein Wechsel der Art der Mitgliedschaft} ist jederzeit
möglich, und wird durch Willenserklärung in Textform gegenüber dem
Vorstand vollzogen.
\strut\end{minipage}\tabularnewline
\bottomrule
\end{longtable}

In der Diskussion wird als Argument angeführt, dass die Formulierung
"`jederzeit"' Missbrauchspotenzial bietet. Beispielsweise könnten
Fördermitglieder kurz vor Mitgliederversammlungen ihre Mitgliedschaft zum
ordentlicher Mitglied umwandeln und somit Stimmrecht erhalten, ohne das ganze
Jahr über den vollen Mitgliedsbeitrag zahlen zu müssen. Es wird daher empfohlen,
stattdessen Formulierungen wie "`zum Anfang des Monats"' etc. einzufügen.

\newpage
Laut dem Schatzmeister wird der Wechsel der Mitgliedschaft bisher zwar nicht mit
diesem Hintergrund ausgenutzt, und der Vorstand kann notfalls auch jederzeit
ohne Angabe von Gründen Mitglieder ausschließen.

\vote{Wechsel der Art der Mitgliedschaft vereinfachen}{3}{9}{10}
Die Abstimmung erfolgt per Handzeichen und ergibt 3 pro, 9 contra, 10
Enthaltungen. Der Änderungsantrag ist somit abgewiesen.


\subsection{Änderungsantrag 5: Übertragung der Mitgliederrechte auf
Repräsentanten}
Die Übertragung des Stimmrechts von Mitgliedern kam auf vergangenen
Versammlungen vor und wurde dort informell behandelt, wird sich aber
auch in Zukunft nicht vermeiden lassen. Daher sollte die Satzung diesen
Fall vorsehen.

Anmerkungen:

\begin{itemize}
\item
  Repräsentant ist wie Mitglied zu behandeln, zählt also insbesondere
  nicht als Gast und hat damit Anwesenheitsrecht auf der
  Mitgliederversammlung. Zusätzlich aus Verständlichkeitsgründen nochmal
  die Formulierung für Beschlussfähigkeit in Abs. 6.
\item
  Beschränkung auf ein vertretenes Mitglied pro Repräsentant, um
  mögliche Stimmenkumulation auf anwesende Mitglieder zu vermeiden, die
  diese Stimmen in Diskussionen als Druckmittel einsetzen könnten.
\item
  Schriftform (§126 BGB) als brauchbarer Mittelweg, kann auch durch
  elektronische Übermittlung des Dokuments (Scan per E-Mail, Fax)
  gewahrt werden. Unter Textform würde z.B. auch IRC fallen, was nicht
  fälschungssicher genug ist, mündliche Aussagen sind nicht ausreichend
  nachvollziehbar.
\end{itemize}

\subsubsection*{§7 (Mitgliederversammlung), Abs. 6}
\begin{longtable}[c]{@{}ll@{}}
\toprule
\begin{minipage}[b]{0.47\columnwidth}\raggedright\strut
alter Text:
\strut\end{minipage} &
\begin{minipage}[b]{0.47\columnwidth}\raggedright\strut
neuer Text:
\strut\end{minipage}\tabularnewline
\midrule
\endhead
\begin{minipage}[t]{0.47\columnwidth}\raggedright\strut
Jede satzungsmäßig einberufene Mitgliederversammlung wird als
beschlussfähig anerkannt, sofern mindestens 23\% der Mitglieder anwesend
sind. Falls dieser geforderte Anteil nicht erreicht wird, ist die darauf
folgende Mitgliederversammlung unabhängig von der Anzahl der erschienen
Mitglieder beschlussfähig. Auf diesen Umstand muss in der Einladung zur
Mitgliederversammlung besonders hingewiesen werden. Jedes ordentliche
Mitglied hat eine Stimme. Fördermitglieder sind berechtigt, an den
Versammlungen ohne Stimmrecht teilzunehmen.
\strut\end{minipage} &
\begin{minipage}[t]{0.47\columnwidth}\raggedright\strut
Jede satzungsmäßig einberufene Mitgliederversammlung wird als
beschlussfähig anerkannt, sofern mindestens 23\% der Mitglieder anwesend
sind \emph{oder durch Repräsentanten repräsentiert werden}. Falls dieser
geforderte Anteil nicht erreicht wird, ist die darauf folgende
Mitgliederversammlung unabhängig von der Anzahl der erschienen
Mitglieder beschlussfähig. Auf diesen Umstand muss in der Einladung zur
Mitgliederversammlung besonders hingewiesen werden. Jedes ordentliche
Mitglied hat eine Stimme. Fördermitglieder sind berechtigt, an den
Versammlungen ohne Stimmrecht teilzunehmen.
\strut\end{minipage}\tabularnewline
\bottomrule
\end{longtable}

\newpage
\subsubsection*{§7 (Mitgliederversammlung), Abs. 10}
\begin{longtable}[c]{@{}ll@{}}
\toprule
\begin{minipage}[b]{0.47\columnwidth}\raggedright\strut
alter Text:
\strut\end{minipage} &
\begin{minipage}[b]{0.47\columnwidth}\raggedright\strut
neuer Text:
\strut\end{minipage}\tabularnewline
\midrule
\endhead
\begin{minipage}[t]{0.47\columnwidth}\raggedright\strut
\emph{(Absatz wird neu eingefügt)}
\strut\end{minipage} &
\begin{minipage}[t]{0.47\columnwidth}\raggedright\strut
\emph{Bei Nichtanwesenheit auf der Mitgliederversammlung kann ein Mitglied
seine Mitgliederrechte für die Dauer der Versammlung auf eine andere
Person (Repräsentant) übertragen. Die Repräsentanten üben die
Mitgliederversammlung betreffend in allen Belangen die Rechte und
Pflichten des Mitglieds aus, das durch sie vertreten wird, und sind
dementsprechend im Rahmen der Mitgliederversammlung wie Mitglieder im
Sinne dieser Satzung zu behandeln. Dabei darf ein Repräsentant zu keiner
Zeit mehr als ein Mitglied vertreten. Falls ein Mitglied als
Repräsentant benannt wird, werden dessen Rechte als Mitglied dadurch
nicht eingeschränkt, in diesem Fall kann ein Mitglied also bis zu zwei
Mitglieder, inkl. sich selbst, repräsentieren. Die Übertragung der
Mitgliederrechte ist dem Vorstand vom entsprechenden Mitglied spätestens
bis zu Beginn der Versammlung in Schriftform bekannt zu machen, wobei
der Repräsentant namentlich benannt werden muss.}
\strut\end{minipage}\tabularnewline
\bottomrule
\end{longtable}

\vote{Übertragung der Mitglieder\-rechte auf Re\-präsentanten}{1}{13}{10}
Dieser Vorschlag wurde im Vorfeld der Versammlung schon auf der Mailingliste
diskutiert und fand dort geteilte Meinungen. Bei der Abstimmung per Handzeichen
gibt es jedoch nur 1 Pro-Stimme, 13 dagegen, 10 Enthaltungen. Der Antrag ist
somit abgelehnt.


%%%%%%%%%%%
%% TOP 4 %%
%%%%%%%%%%%
\section{Optionsrecht Mietvertrag}
Der Mietvertrag für den Space 2.0 besagt in §2.4: \emph{"`Der Mieter kann von
der Ausübung des Optionsrechts nur schriftlich bis spätestens 6 Monate vor
Vertragsablauf bis zum 31.03.2016 Gebrauch machen."'} Über eine Verlängerung des
Mietvertrages bis effektiv 31.9.2019 soll von der Mitgliederversammlung
abgestimmt werden. Bei den laufenden Kosten von 10.000€/Jahr handelt es sich
also um eine Abstimmung über insgesamt 30.000€.

In der Diskussion werden folgende Punkte genannt:
\begin{itemize}
  \item \textbf{Vorteile der aktuellen Räumlichkeit:}
  \begin{itemize}
    \item günstige Lage: innenstadtnah, uninah, gute Parkplatzsituation,
      Gewerbegebiet lässt Lärm zu
    \item nach unserem Bedarf ausgebaut
  \end{itemize}

  \item \textbf{Nachteile der aktuellen Räumlichkeit:}
  \begin{itemize}
    \item nicht barrierefrei
    \item Platz auch hier begrenzt, Wachstum hält Schritt
  \end{itemize}

  \item \textbf{Sonstige Punkte:}
  \begin{itemize}
    \item Platzangebot ist auch noch durch mehr Ordnung verbesserbar, Probleme
      im Moment hauptsächlich bei gleichzeitig stattfindenden Terminen
    \item Umzug erfordert aufmerksame Immobiliensuche bei aktuell schwieriger
      Marktlage
    \item nach aktueller Finanzlage nur etwa 100€/Monat mehr investierbar
    \item sonst: Beiträge erhöhen, Pledgefunding betreiben
    \item wir haben jetzt Übung im Renovieren
    \item möglicherweise will Vermieterin nachverhandeln
    \item Untervermietung für 5€/m² ist im Notfall auch immer noch möglich
  \end{itemize}
\end{itemize}

Es spricht also anscheinend vieles dafür, in den aktuellen Räumlichkeiten zu
bleiben und den Mietvertrag zu verlängern.

\consensus{Die Mitgliederversammlung beauftragt den Vorstand, das Optionsrecht
des aktuellen Mietvertrags auf die Verlängerung um weitere drei Jahre zu nutzen.}
\paragraph{Beschlussvorlage:} \emph{Die Mitgliederversammlung beauftragt den
Vorstand, das Optionsrecht des aktuellen Mietvertrags auf die Verlängerung um
weitere drei Jahre zu nutzen.}

Diese Beschlussvorlage wird per Handzeichen einstimmig angenommen.

%%%%%%%%%%%
%% TOP 5 %%
%%%%%%%%%%%
\section{Sonstiges}

\subsection{Ansprechpartner für Konflikte} Nach den Diskussionen um die
Gestaltung des EasterHeggs hatte der Vorstand sich in §5 seiner Geschäftsordnung
zusätzlich die Aufgabe gegeben, als Ansprechpartner für Konflikte
bereitzustehen.\footnote{siehe
\url{https://stratum0.org/wiki/Vorstand/Geschäftsordnung}} Dies soll auch
weiterhin gegeben sein. Es sollte aber zusätzlich eine Vertrauensperson außerhalb
des Vorstands zur Verfügung stehen, falls Konflikte mit dem Vorstand auftreten
oder Leute sich aus anderen Gründen nicht an den Vorstand wenden wollen. Falls
es eine Person aus dem Verein gibt, die diese Funktion erfüllen würde, sollte
die Mitgliederversammlung diese Person damit beauftragen. Die entsprechenden
Prozesse sind im Vorstand schon definiert worden und könnten der Person als
Vorlage dienen.

Der Vorschlag, diesem Posten das Recht zu Anwesenheit bei
den nichtöffentlichen Teilen von Vorstandssitzungen zu geben, wird vom Vorstand
abgelehnt, da in den nichtöffentlichen Teilen oft auch andere
datenschutzrechtliche Dinge behandelt werden. Außerdem könnte dieses Recht
kontraproduktiv wirken, falls sich Leute mit Konflikten explizit an den Vorstand
wenden, weil sie nicht mit der gewählten Vertrauensperson sprechen wollen.

\consensus{Die Mitgliederversammlung benennt eine Entität als
zusätzlichen Ansprechpartner, die neben dem Vorstand als Vertrauensperson im
Verein zur Verfügung steht.}
\paragraph{Beschlussvorlage:} \emph{Die Mitgliederversammlung benennt eine Entität als
zusätzlichen Ansprechpartner, die neben dem Vorstand als Vertrauensperson im
Verein zur Verfügung steht.}

Die Beschlussvorlage wird per Handzeichen abgestimmt und einstimmig ohne
Enthaltungen angenommen.

\consensus{dStulle als Vertrauensperson beauftragen}
dStulle meldet sich freiwillig für diesen Posten. Es wird per Handzeichen
darüber abgestimmt, dStulle mit dieser Funktion zu beauftragen. Die Abstimmung
fällt mehrheitlich dafür aus. Seine Kontaktdaten sind im Wiki\footnote{siehe
\url{https://stratum0.org/wiki/Benutzer:dStulle}} ersichtlich.

\subsection{Belüftung im Frickelraum}
Vor die halbhohen Fenster im Chillraum und im Frickelraum wurden wegen besserer
Isolierung neue Scheiben gesetzt, die sich nicht mehr öffnen lassen. Da die
Luftzirkulation im Frickelraum somit nicht mehr stattfinden kann, sollte dort
ein Lüfter eingebaut werden. Der Architekt hatte dafür einen Preis von etwa 500€
genannt, die Vermieterin würde sich zur Hälfte an den Kosten beteiligen.
Zusätzlich gab es schon unsererseits beim Ausbau im Rahmen des Space 2.0
Planungen für eine Belüftungsanlage, die sich auf etwa 350€ beliefen.

\consensus{Belüftung im Frickelraum einbauen}
Generell ist die Versammlung dafür, den Frickelraum besser zu belüften. Der
Vorstand wird sich dementsprechend weiter darum kümmern.

\subsection{Sport als Vereinszweck}
Es wird eine kurze, generelle Diskussion angestoßen, ob Sport in den
Vereinszweck aufgenommen werden soll. Damit könnten z.~B. CTFs oder Lockpicking
abgedeckt werden. Möglicherweise wäre es somit möglich, Geld z.~B. beim
Stadtsportbund zu beantragen.

Es ist noch nicht klar, inwiefern sich der Vereinszweck Sport auf die
Steuerbegünstigung auswirkt. Da diese Frage auch schon in einer Vorstandssitzung
auftrat, sollen in Zukunft mehr Informationen dazu zusammengetragen werden.

\subsection{Termin Nächste Mitgliederversammlung}
Der Termin dieser Mitgliederversammlung wurde relativ kurzfristig angekündigt,
weshalb einige Mitglieder nicht erscheinen konnten. Deshalb wird an dieser
Stelle schon einmal angekündigt, dass der erste Sonntag im Dezember, also der 4.
Dezember 2016 als Termin für die nächste Mitgliederversammlung ins Auge gefasst
wird. Der Termin scheint bisher allen Anwesenden zu passen.

\begin{description}
  \item[Veranstaltung geschlossen] um 18:58
\end{description}


%%%%%%%%%%%%%%%%%%%%
%% Unterschriften %%
%%%%%%%%%%%%%%%%%%%%
\cleardoublepage
\section*{Unterschriften}
\vspace{0.7cm}
\noindent Protokollführer: \hrulefill\hfill\phantom{c}\par
\vspace{0.7cm}
\noindent Vorstandsvorsitzender: \hrulefill\hfill\phantom{c}\par
\vspace{0.7cm}
\noindent Stellv. Vorsitzender: \hrulefill\hfill\phantom{c}\par
\vspace{0.7cm}
\noindent Schatzmeister: \hrulefill\hfill\phantom{c}\par
\vspace{0.7cm}
\noindent Beisitzer: \hrulefill\hfill\phantom{c}\par
\vspace{0.7cm}
\noindent Beisitzerin: \hrulefill\hfill\phantom{c}\par
\vspace{0.7cm}
\noindent Beisitzer: \hrulefill\hfill\phantom{c}\par

\end{document}
